\begin{proposicion}[Convergencia uniforme]\label{prop:conv_unif}
Cuando $x\downarrow0$, $r_{R,x},r_{C,x},u_{D,x},u_{A,x}$ convergen uniformemente a versiones continuas
y positivas de $r_{R,0},r_{C,0},u_{D,0},u_{A,0}$, que cumplen las cotas de la
Proposición~\ref{prop:cotas} y las ecuaciones~\eqref{eq:filas}--\eqref{eq:diagonales} con $x=0$ y $b=b_0$.
\end{proposicion}

\begin{proof}
Escribimos las integrales de~\eqref{eq:filas}--\eqref{eq:diagonales} como $\Phi_f$, de modo que
$r_R=1/(x+b\Phi_R)$, $r_C=1/(x+b\Phi_C)$, $u_D=1/(1+b\Phi_D)$ y $u_A=1/(1+b\Phi_A)$, y sea $\Phi_{f,0}$
la misma integral formada con los límites. En cada integral de fibra, todo factor ajeno a la familia
de la fibra se evalúa en una coordenada afín de pendiente de módulo uno o dos; no hay pendiente cero,
porque las cuatro familias no son paralelas entre sí. La diferencia de productos se acota por la suma
de las diferencias de sus factores, acotados por $M$, y un cambio de variable da
\begin{equation}\label{eq:fibras_unif}
 \sup_{t\in I_f}|\Phi_f(t)-\Phi_{f,0}(t)|\le C\sum_{g\ne f}\|y_{g,x}-y_{g,0}\|_{L^1(I_g)},
\end{equation}
donde $y_g$ es el factor normalizado de la familia $g$ ($r_R$, $r_C$, $u_D$ o $u_A$) y $C$ depende sólo
de $M$. Esta estimación vale antes de saber que los límites son continuos. Las integrales límite
$\Phi_{f,0}$ son continuas, por la continuidad de las traslaciones en $L^1$ y porque los extremos de las
fibras se mueven continuamente, como en el Lema~\ref{lem:vacancias}. Por la
Proposición~\ref{prop:conv_L2}, $\Phi_f\to\Phi_{f,0}$ uniformemente. Los denominadores están separados
de cero: $b\Phi_R,b\Phi_C\ge\kappa^2/(2(1+b_+))$ y $1+b\Phi_D,1+b\Phi_A\ge1$. Las fórmulas racionales
dan entonces la convergencia uniforme a $1/(b_0\Phi_{R,0})$, etc., que son continuas, positivas e
iguales en casi todo punto a los límites en $L^2$.
\end{proof}

Para el teorema de la función implícita usamos el espacio de Banach
$X=C(I_R)\oplus C(I_C)\oplus C(I_D)\oplus C(I_A)$, con incógnitas
$z=(z_R,z_C,z_D,z_A)=(\log r_R,\log r_C,\log u_D,\log u_A)\in X$ y $\gamma=\log b$. Para
$(x,z,\gamma)\in\mathbb R\times X\times\mathbb R$ sean $\nu=b\,r_Rr_Cu_Du_A$ y
\[
 E_R=xr_R+T_R\nu-1,\quad E_C=xr_C+T_C\nu-1,\quad E_D=u_D+T_D\nu-1,\quad E_A=u_A+T_A\nu-1,
\]
\[
 n_R=\int r_R-1,\qquad n_C=\int r_C-1 .
\]
Son aplicaciones analíticas reales entre espacios de Banach (exponenciales, productos, proyecciones
lineales acotadas e integrales). Como $\int T_R\nu=\int T_C\nu$, hay una redundancia exacta:
$\int E_R-\int E_C=x(n_R-n_C)$. Sean $Z\subset X$ el subespacio cerrado $\{\int e_R=\int e_C\}$, de
codimensión uno, y $P(e)=(e_R,\,e_C-\int e_C+\int e_R,\,e_D,\,e_A)$, que lleva $X$ sobre $Z$ y es la
identidad en $Z$. Consideramos
\begin{equation}\label{eq:sistema}
 \mathscr E(x,z,\gamma)=(PE,n_R,n_C)=0\qquad\text{en }Z\oplus\mathbb R^2 .
\end{equation}
Si $PE=0$ y $n_R=n_C=0$, la redundancia da $\int E_C=\int E_R=0$ y luego $E=0$, también para
$x\ne0$; así~\eqref{eq:sistema} equivale a $E=0$ junto con las dos normalizaciones. Por el
Lema~\ref{lem:vacancias} y la Proposición~\ref{prop:conv_unif}, los maximizadores ($0<x\le\frac12$) y
el límite ($x=0$) lo satisfacen.

\begin{teorema}[Rama analítica en el tablero lleno]\label{thm:rama}
Existen $\varepsilon>0$ y una aplicación analítica real $x\mapsto(z(x),\gamma(x))$ de
$(-\varepsilon,\varepsilon)$ en $X\times\mathbb R$ que resuelve~\eqref{eq:sistema}, con
\[
 (z(0),\gamma(0))=(\log r_{R,0},\log r_{C,0},\log u_{D,0},\log u_{A,0},\log b_0),
\]
y que para $0<x<\varepsilon$ coincide con el perfil óptimo:
\[
 (z(x),\gamma(x))=(\log r_{R,x},\log r_{C,x},\log u_{D,x},\log u_{A,x},\log b_x).
\]
\end{teorema}

\begin{proof}
Para una variación $(h,\beta)=(h_R,h_C,h_D,h_A,\beta)\in X\oplus\mathbb R$ ponemos
$w=\beta+\sum_fh_f\circ\ell_f$, de modo que $\dot\nu=\nu_1w$ en $x=0$. La parte de la derivada de
$E_f$ que actúa sobre $h_f$ es $h_f$: para filas y columnas porque $T_R\nu_1=T_C\nu_1=1$, y en una
diagonal porque $(u_{D,0}+m_D[\nu_1])h_D=h_D$. Los operadores cruzados entre familias distintas son
compactos por Arzelà--Ascoli: parametrizados por la coordenada de la otra familia, tienen núcleos
continuos en el dominio que permite la fibra (extenderlos por cero puede crear saltos en el borde de
ese dominio, lo que no afecta al argumento); para dos parámetros de fibra próximos se usa la
continuidad uniforme del núcleo en el dominio común, y la diferencia simétrica de los intervalos de
integración tiene longitud $O(|t-t'|)$. Para $D\leftrightarrow A$ el cambio de variable tiene
jacobiano $\frac12$ y extremos $[|d|,2-|d|]$. El término en $\beta$ tiene rango uno. Así, la derivada
de $E$ es un operador de Fredholm de $X\oplus\mathbb R$ en $X$ de índice uno. En $x=0$ su imagen está
en $Z$, porque la derivada de la redundancia se anula, lo que eleva el índice a dos como operador con
valores en $Z$; las dos ecuaciones de normalización lo devuelven a cero.

Si $(h,\beta)$ está en el núcleo de la derivada de~\eqref{eq:sistema} en $x=0$, entonces
\begin{gather*}
 T_R(\nu_1w)=T_C(\nu_1w)=0,\qquad T_D(\nu_1w)=-u_{D,0}h_D,\qquad T_A(\nu_1w)=-u_{A,0}h_A,\\
 \int r_{R,0}h_R=\int r_{C,0}h_C=0 .
\end{gather*}
Multiplicando por las componentes de $w$, integrando y usando la ecuación de filas para anular
$\beta\int\nu_1w=\beta\int T_R(\nu_1w)=0$, se obtiene
\begin{equation}\label{eq:energia}
 \int_{[0,1]^2}\nu_1w^2+\int_{I_D}u_{D,0}h_D^2+\int_{I_A}u_{A,0}h_A^2=0 .
\end{equation}
Por positividad, $h_D=h_A=0$ y $\beta+h_R(r)+h_C(c)=0$ en el cuadrado; luego $h_R$ y $h_C$ son
constantes, las normalizaciones las anulan y entonces $\beta=0$. El núcleo es trivial y, con índice
cero, la derivada es un isomorfismo con inversa acotada. El teorema analítico de la función
implícita~\cite{Zeidler} da un entorno $U$ de $(0,z(0),\gamma(0))$ y una rama analítica cuyos puntos
son exactamente las soluciones de~\eqref{eq:sistema} en $U$. Por la Proposición~\ref{prop:conv_unif} y las
cotas inferiores, los logaritmos de los perfiles óptimos convergen uniformemente cuando $x\downarrow0$;
para $x>0$ pequeño están en $U$ y coinciden con la rama.
\end{proof}

Las Proposiciones~\ref{prop:cotas}, \ref{prop:conv_L2} y~\ref{prop:conv_unif} y el
Teorema~\ref{thm:rama} prueban el Teorema~\ref{thm:lleno}.
